% =============================================================================
% Preámbulo compartido. Sigue el estilo de project_document/ para que los dos
% documentos del repositorio se lean como uno.
% =============================================================================

\documentclass[11pt,a4paper]{article}

\usepackage[utf8]{inputenc}
\usepackage[T1]{fontenc}
\usepackage[spanish,es-noquoting]{babel}
\usepackage[margin=2.4cm]{geometry}
\usepackage{microtype}
\usepackage{parskip}
\usepackage{enumitem}
\usepackage{xcolor}
\usepackage{booktabs}
\usepackage{longtable}
\usepackage{tabularx}
\usepackage{xltabular}
\usepackage{array}
\usepackage{amsmath}
\usepackage{amssymb}
\usepackage{titlesec}
\usepackage{fancyhdr}
\usepackage{caption}
\usepackage[hidelinks]{hyperref}

% --- color -------------------------------------------------------------------
\definecolor{acc}{HTML}{146B57}
\definecolor{dim}{HTML}{5C6B65}
\definecolor{rule}{HTML}{D8DEDA}
\definecolor{aviso}{HTML}{8A6410}

\hypersetup{
  unicode=true,
  colorlinks=true,
  linkcolor=acc,
  urlcolor=acc,
  citecolor=acc,
  pdftitle={PROTEA. Diseño de la campaña de experimentación},
  pdfauthor={Francisco Miguel Pérez Canales}
}

% --- títulos -----------------------------------------------------------------
\titleformat{\section}{\Large\bfseries\color{acc}}{\thesection}{0.8em}{}
\titleformat{\subsection}{\large\bfseries}{\thesubsection}{0.8em}{}
\titleformat{\subsubsection}{\normalsize\bfseries}{\thesubsubsection}{0.7em}{}
\titlespacing*{\section}{0pt}{2.2ex plus 1ex minus .2ex}{1.2ex plus .2ex}

% --- cabeceras ---------------------------------------------------------------
\pagestyle{fancy}
\fancyhf{}
\fancyhead[L]{\footnotesize\color{dim} PROTEA. Diseño de la campaña}
\fancyhead[R]{\footnotesize\color{dim}\thepage}
\renewcommand{\headrulewidth}{0.4pt}
\renewcommand{\headrule}{\hbox to\headwidth{\color{rule}\leaders\hrule height \headrulewidth\hfill}}

% --- tablas ------------------------------------------------------------------
\captionsetup{font=small,labelfont={bf,color=acc},skip=6pt}
\renewcommand{\arraystretch}{1.18}
\newcolumntype{L}[1]{>{\raggedright\arraybackslash}p{#1}}
\newcolumntype{C}[1]{>{\centering\arraybackslash}p{#1}}
\newcolumntype{R}[1]{>{\raggedleft\arraybackslash}p{#1}}

% --- el entorno de definiciones ---------------------------------------------
% Toda la sección 2 pasa por aquí. Un término definido aquí es el único
% significado que tiene en el resto del documento.
\newcounter{defn}
\makeatletter
% \term{Nombre}{cuerpo}  o  \term[etiqueta]{Nombre}{cuerpo}
% La etiqueta permite citar una definición con \ref, que imprime su número.
\newcommand{\term}[3][]{%
  \refstepcounter{defn}%
  \edef\pr@tmp{#1}%
  \ifx\pr@tmp\@empty\else\label{#1}\fi
  \par\vspace{0.65em}%
  \noindent\textbf{D\thedefn. #2.}\enspace #3\par\vspace{0.15em}}
\makeatother

% Un término del código. Se deja en inglés a propósito, véase la sección 2.0.
\newcommand{\cod}[1]{\texttt{\small #1}}

% Una cifra medida, con su procedencia.
\newcommand{\medido}[1]{\textcolor{dim}{\footnotesize (medido: #1)}}

% Un aviso que exige decision del investigador.
\newenvironment{aviso}[1]{%
  \par\vspace{0.9em}\noindent
  \begingroup\color{aviso}\rule{\linewidth}{1pt}\endgroup\par\vspace{0.35em}%
  \noindent\textbf{\color{aviso}#1}\par\vspace{0.25em}\noindent\ignorespaces}%
 {\par\vspace{0.35em}\noindent
  \begingroup\color{aviso}\rule{\linewidth}{1pt}\endgroup\par\vspace{0.9em}}


\title{\textbf{PROTEA}\\[0.35em]
  \large Diseño de la campaña de experimentación\\
  \normalsize Dimensiones, ejes, parámetros y regla de cierre\\[0.7em]
  \small Informe para el centro de investigación}

\author{Francisco Miguel Pérez Canales\\
  \small Universidad de Sevilla\\
  \small \texttt{frapercan1@alum.us.es}}

\date{Septiembre de 2026}

\begin{document}
\maketitle
\thispagestyle{empty}

\begin{abstract}
\noindent
Este documento define la campaña de experimentación de PROTEA: qué se decide,
sobre qué dimensiones, con qué parámetros, y bajo qué regla una diferencia
medida cuenta como una decisión tomada. No presenta resultados. Su objeto es la
estructura, y el criterio de éxito es que ningún término quede sin definir y
ningún parámetro sin declarar.

El diseño está pensado para ser invariante a la escala. La campaña que se
ejecuta hoy está limitada por la potencia disponible en una máquina de cómputo,
y la campaña que el centro podría ejecutar sobre cinco máquinas es la misma
campaña con una rejilla más ancha. La sección~\ref{sec:escalas} separa
explícitamente lo que cambia con la escala, que es el tamaño de la rejilla y el
número de niveles que se pueden comparar a la vez, de lo que no cambia, que son
los diez nodos de decisión, las definiciones, la unidad de reporte y la regla
de cierre. Ese reparto es el resultado principal del documento: más cómputo
compra anchura de rejilla, no un método distinto.

Las cifras que aparecen son medidas sobre el registro actual y llevan su
procedencia. Donde una propiedad no se ha establecido, se dice.
\end{abstract}

\vspace{1em}
\begin{center}
\begin{tabularx}{0.86\linewidth}{@{}lX@{}}
\toprule
\multicolumn{2}{@{}l}{\textbf{Resumen de la estructura}}\\
\midrule
Decisiones      & 10 nodos tipados, en 10 etapas ordenadas\\
Firmeza         & 5 palabras, en un orden de pruebas fijo\\
Unidad de reporte & 9 paneles obligatorios, refinables sobre 7 ejes de estrato\\
Marco           & 6 campos sellados hoy, 8 propuestos\\
Parámetros      & 15 en predicción, 7 en evaluación, 11 en el sustrato\\
Estadístico     & bootstrap pareado por proteína, IC del 95\,\%\\
Regla de cierre & unanimidad sobre las celdas con poder, o conjunto supervivente\\
\bottomrule
\end{tabularx}
\end{center}

\clearpage
\tableofcontents
\clearpage

\section{Propósito y alcance}
\label{sec:proposito}

\subsection{La pregunta}

La campaña responde a una sola pregunta, y está formulada de modo que su
respuesta puede ser negativa:

\begin{quote}
\emph{Cuáles de las decisiones que definen un sistema de transferencia de
anotación basado en modelos de lenguaje y vecinos más cercanos están separadas
por evidencia sobre una ventana temporal sin fuga, y cuáles están simplemente
elegidas.}
\end{quote}

La distinción entre \emph{separada por evidencia} y \emph{elegida} no es
retórica. Está implementada: cada decisión del sistema recibe exactamente una
de cinco palabras, y la más fuerte de las cinco solo es alcanzable cuando se
declaró un umbral de comparación antes de medir y la comparación lo superó. La
sección~\ref{sec:grafo} define las cinco palabras y el orden en que se prueban.

De esto se sigue el compromiso central del diseño: la campaña se compromete por
adelantado a publicar como \emph{elegida} cualquier decisión que no alcance la
palabra fuerte, en lugar de presentarla como un hallazgo. Una campaña que solo
puede producir hallazgos no está midiendo nada.

\subsection{Lo que este documento no es}

No es un informe de resultados. Las cifras que contiene son de tres clases y
están marcadas como tales: poblaciones y costes medidos sobre el registro
actual, que fijan la escala; cotas de resolución, que fijan lo que es
observable; y valores de parámetro instanciados, que fijan la rejilla. No hay
ninguna afirmación comparativa entre configuraciones.

Tampoco es un plan de trabajo cerrado. Las decisiones que corresponden al
investigador están recogidas en la sección~\ref{sec:procedimiento} como puertas
que pueden fallar, con la entrada sobre la que fallan nombrada en cada caso.

\subsection{Por qué el diseño se escribe antes de los resultados}

Porque el defecto que esta campaña existe para corregir no es un error de
cálculo. Es estructural, y se enuncia en una frase: \textbf{un proyecto de esta
forma escribe sus decisiones en prosa y sus datos en tablas, y nada conecta las
dos}. Cuando eso ocurre, una prohibición escrita en un documento no impide la
comparación que prohíbe, un umbral declarado no puede negarse, y una fila
retirada sigue gobernando la lectura de un resultado porque la consulta que la
lee no menciona su estado.

La consecuencia práctica es que \emph{declarar} algo tiene que costar más que
escribirlo. En este diseño una declaración vale si y solo si tiene tres partes,
y no vale nada con dos:

\begin{enumerate}[leftmargin=1.5em,itemsep=0.25em]
  \item \textbf{Una fila}, en una columna tipada, direccionable por el código
        que debe obedecerla. Prosa en un fichero de texto no es una
        declaración; tampoco lo es una clave libre en un blob que ninguna
        consulta nombra.
  \item \textbf{Un guardia que pueda negarse}, es decir, algún camino que lea
        esa fila y \emph{levante una excepción} antes de que el número se
        produzca o la comparación se responda. Un guardia que registra un aviso,
        o que devuelve un valor por omisión, no es un guardia.
  \item \textbf{Dos pruebas}: una de que se niega cuando la regla se viola, y
        otra de que actúa cuando la regla se respeta. Las dos mitades son
        necesarias. Un guardia probado solo en la negativa puede ser una
        función que siempre levanta; probado solo en el camino feliz, puede ser
        una que nunca dispara.
\end{enumerate}

Todo cambio de instrumentación que propone la sección~\ref{sec:procedimiento}
se ajusta a esa plantilla, y la razón por la que este documento define cada
término es la misma: un guardia solo puede negarse a algo que esté definido.

\subsection{Audiencia y uso}

El documento está escrito para ser leído por un grupo que no ha seguido el
desarrollo. Por eso la sección~\ref{sec:definiciones} define todo el
vocabulario antes de usarlo, incluidos los términos que parecen obvios. Los
nombres técnicos se dejan en inglés, en \texttt{\small esta tipografía}, porque son los
identificadores que existen en el código y en el esquema de la base de datos.
Traducirlos crearía un segundo vocabulario que habría que mantener sincronizado
con el primero, y un conjunto cerrado escrito dos veces es la forma más común de
que las dos copias dejen de coincidir sin que nadie lo note.

\section{Definiciones}
\label{sec:definiciones}

Esta sección es normativa. Un término definido aquí tiene ese significado y
ningún otro en el resto del documento, y cada definición nombra el sitio del
código o del esquema donde el objeto existe de verdad. Las definiciones están
numeradas para poder citarlas.

\subsection{Nota sobre el vocabulario}

Los nombres técnicos se conservan en inglés. No es un descuido de redacción: son
los identificadores del código y de las columnas, y el documento se usa para
discutir sobre ellos. Donde el castellano aporta claridad se usa la traducción
como glosa, nunca como sustituto.

\subsection{Los objetos del registro}

\term{Protein}{Una secuencia de aminoácidos con un accession de UniProt. Es la
entidad sobre la que se predice y la unidad sobre la que se estratifica.}

\term{GO term}{Un nodo de la Gene Ontology, identificado por su \cod{go\_id}
textual. Pertenece a exactamente uno de los tres aspectos (D\ref{def:aspect}).
Un término existe siempre dentro de un \emph{ontology snapshot}: el mismo
\cod{go\_id} es una fila distinta en dos snapshots, y confundir las dos es la
forma habitual de comparar dos grafos distintos creyendo que es uno.}

\term{Annotation}{La afirmación de que una proteína tiene una función, con un
código de evidencia que dice cómo se supo. Vive en \cod{protein\_go\_annotation}
y pertenece a un \emph{annotation set}.}

\term{Annotation set}{El contenido de una release concreta de GOA, cargada
completa. Es el objeto que la campaña fecha: una release nombra un fichero y no
fecha nada por sí misma, de modo que la fecha de publicación se almacena aparte.}

\term{Evidence code}{La etiqueta de tres letras que clasifica el soporte de una
anotación. La campaña distingue dos grupos que importan para todo lo demás: los
códigos \emph{experimentales} y de \emph{alto rendimiento} (\cod{EXP}, \cod{IDA},
\cod{IPI}, \cod{IMP}, \cod{IGI}, \cod{IEP}, \cod{HTP}, \cod{HDA}, \cod{HMP},
\cod{HGI}, \cod{HEP}, más \cod{IC} y \cod{TAS}), y el resto, que es
mayoritariamente inferencia computacional (\cod{IEA}). La distinción es la que
usan CAFA y LAFA para construir su verdad de referencia, y aquí se usa en dos
sitios distintos que conviene no confundir: decide qué anotaciones pueden
\emph{donar} (D\ref{def:bank}) y decide qué anotaciones cuentan como verdad al
\emph{evaluar}.}

\term{Query}{Una proteína para la que se pide predicción en una corrida. El
conjunto de queries de la campaña es un \cod{query\_set} fijo.}

\term{Donor}{Una proteína del banco cuya anotación se transfiere a una query
porque su vector está cerca. Es el mecanismo entero del método base: no se
predice una función, se copia la de un vecino.}

\term{Candidate}{Un par (query, GO term) propuesto por el sistema, con una
puntuación. La columna \cod{go\_prediction.ref\_protein\_accession} es
\cod{NOT NULL}, de modo que el esquema prohíbe un candidato sin donante en lugar
de limitarse a registrar que nadie escribió uno.}

\term{Arm}{Una configuración completa, ejecutada de principio a fin, que produce
un conjunto de predicciones evaluable. Es la unidad que se compara con otra. Dos
arms difieren en al menos un parámetro, y nombrar cuál es precisamente el
problema que resuelve la noción de \emph{nivel} (D\ref{def:nivel}).}

\term{Prediction set}{El producto de un arm: el conjunto de candidatos con su
procedencia. Lleva en \cod{meta} los parámetros con los que se generó.}

\term{Evaluation set}{Un par de annotation sets que define una ventana temporal:
el de la cota inferior, del que pueden venir los donantes, y el de la cota
superior, que aporta la verdad de referencia. Nada más es una ventana.}

\term{Evaluation result}{Una fila con las métricas de un prediction set contra un
evaluation set. Es el único sitio donde vive un número publicable.}

\subsection{El marco}

\term[def:marco]{Frame}{El conjunto de condiciones bajo las que un número se
lee, y que por tanto deben coincidir para que dos números sean comparables. Hoy
son seis campos: el evaluation set, el pivot snapshot, el conjunto de
information accretion, la etiqueta de ventana temporal, el número máximo de
términos y la distancia máxima. La sección~\ref{sec:marco} argumenta que dos de
ellos están mal elegidos y que faltan dos.}

\term{Window}{El intervalo temporal de un evaluation set, escrito como
\emph{inferior} $\rightarrow$ \emph{superior}. La campaña selecciona sobre
$220 \rightarrow 227$. La cota inferior es la única release de la que un donante
puede venir sin leer el futuro, y de ahí que el frame fije el banco.}

\term{Pivot snapshot}{El ontology snapshot bajo el que se resuelven los términos
al comparar predicción y verdad. Importa porque la ontología cambia entre
releases: un término obsoleto o reubicado cambia de posición en el grafo, y con
ella cambia la propagación a ancestros.}

\term{Information accretion}{El peso por término que usa CAFA para que acertar
un término raro valga más que acertar uno trivialmente frecuente. Se calcula
sobre el corpus de la cota inferior y se congela en un
\cod{information\_accretion\_set}. Es el \emph{w} de \cod{fmax\_w}.}

\term{Frame seal, o frame digest}{El resumen criptográfico de los campos del
frame, almacenado en cada evaluation result. Es el único sello real: dos
resultados son comparables si y solo si su digest coincide. La columna llamada
\cod{frame}, que contiene el texto \cod{internal}, es una etiqueta del arnés de
ejecución y no puede sostener esa garantía.}

\term{Leak, o fuga}{Que un número use, directa o indirectamente, información
posterior a la cota superior de su ventana. La campaña la impide de tres
maneras, en orden de fuerza: por ausencia, no cargando en la base las releases
que no debe mirar; por guardia, negándose a evaluar una codificación cuyo corte
de entrenamiento cae dentro de la ventana; y por protocolo, que es la más débil
y la única que depende de la disciplina de quien opera.}

\subsection{Las diez decisiones}
\label{sec:diez}

El sistema se describe como una cadena de diez decisiones ordenadas. Cada una
es un \emph{nodo} (D\ref{def:nodo}) y tiene una pregunta propia. El orden es el
del flujo de datos: una decisión de etapa $n$ solo tiene sentido si las de etapa
menor ya están tomadas.

\term[def:substrate]{Substrate (etapa 1). El espacio vectorial en el que se
define la palabra \emph{vecino}}{Es la configuración de codificación completa, no
solo el nombre del modelo. Once campos la determinan: \cod{model\_name},
\cod{model\_backend}, \cod{layer\_indices}, \cod{layer\_agg}, \cod{pooling},
\cod{normalize}, \cod{normalize\_residues}, \cod{max\_length},
\cod{use\_chunking}, \cod{chunk\_size} y \cod{chunk\_overlap}. Cambiar cualquiera
cambia las distancias, y con ellas cambia quién es vecino de quién, qué donantes
se recuperan y qué se predice. Por eso el substrate es la primera decisión
sustantiva y no un detalle de implementación.

Es útil decir qué \emph{no} es. No es un clasificador: nada se entrena aquí. Es
solo la función que convierte una secuencia en un vector. Todo lo que el método
base hace después consiste en medir distancias en ese espacio y copiar
anotaciones de los más cercanos.

Nada en el frame nombra una representación, de modo que el substrate nunca está
fijado por la ventana: si solo hay un nivel instanciado, es un valor que nadie
eligió y no uno que el marco impuso.}

\term[def:bank]{Bank (etapa 2). Quien tiene derecho a donar, y qué se le permite
decir}{Es la decisión que el resto del documento llama \emph{banco}. Se compone
de dos mitades que no pueden decidirse por separado, y por eso forman un solo
nodo:

\begin{itemize}[leftmargin=1.4em,itemsep=0.2em]
  \item El \textbf{corpus}: de qué annotation set salen los donantes. Lo fija el
        frame, porque la cota inferior de la ventana es la única release de la
        que se puede donar sin leer el futuro.
  \item La \textbf{política de donantes}: qué subconjunto de ese corpus puede
        donar de hecho. Son \cod{donor\_reviewed\_only}, la lista
        \cod{donor\_evidence\_codes} y la lista \cod{donor\_exclusions} de
        prefijos de referencia excluidos.
\end{itemize}

\begin{aviso}{Tres interruptores que la prosa atribuye al banco y el código al retriever}
\cod{exclude\_self\_neighbour}, \cod{expand\_votes\_to\_ancestors} y
\cod{aspect\_separated\_knn} modifican el acto de donar y es natural leerlos como
política de banco. En el árbol no lo son: \cod{\_BANK\_FIELDS} contiene solo
\cod{bank\_source}, \cod{bank\_version}, \cod{donor\_reviewed\_only},
\cod{donor\_evidence\_codes} y \cod{donor\_exclusions}, mientras los tres booleanos
están en \cod{\_RETRIEVER\_FIELDS} bajo un comentario que dice literalmente
\emph{quién puede ser vecino, que es este nodo y no el banco}.

\textbf{La consecuencia hay que decirla entera.} El eje que la campaña declaró
cerrado como \emph{banco y donante} se midió sobre esos tres booleanos más la
política de donantes. Los tres booleanos cierran, en el código, el nodo
\cod{retriever}; el nodo \cod{bank} sigue en \cod{unpowered} porque su único
nivel puntuado es el corpus 220. Un lector no puede saber hoy si lo cerrado fue
el banco o el retriever, y el nodo \cod{bank} solo sale de \cod{unpowered}
cargando un segundo corpus, que es exactamente la escalera de recencia. Reconciliar
la agrupación en prosa con la agrupación implementada es una decisión pendiente,
no una errata.
\end{aviso}

Son un nodo porque un corpus y el filtro que se le aplica definen un banco entre
los dos, y porque un grupo de campos inseparables se sostiene con la firmeza de
su miembro más flojo: si el corpus está fijado por el marco pero la política está
vacía, el nodo entero es un valor heredado.

La distinción con el substrate, que es la confusión más fácil de cometer, se dice
en una frase: \textbf{el substrate decide quién está cerca; el bank decide a
quién de los que están cerca se le hace caso}. Un banco más permisivo admite
donantes cuya anotación se apoya en inferencia computacional, lo que aumenta la
cobertura y baja la fiabilidad por donante. Un banco restringido a evidencia
experimental hace lo contrario. Cuál de los dos gana no es evidente a priori y
es precisamente lo que la campaña mide.

Una propiedad del banco que conviene tener presente porque contamina varias
lecturas: bajo el protocolo temporal, la query se recupera a sí misma desde el
banco de la release anterior si ya estaba anotada entonces. Ese \emph{self hit}
está a distancia cero y es el donante más cercano de casi todas las proteínas
con anotación previa. \cod{exclude\_self\_neighbour} es el interruptor que lo
retira, y varias definiciones de la sección~\ref{sec:unidad} se resuelven
explícitamente sobre donantes de secuencia distinta a la query por esta razón.}

\term[def:retriever]{Retriever (etapa 3). Cómo se extraen candidatos del banco, y
a qué profundidad}{Dos parámetros y un cuidado. La \textbf{métrica} de distancia
(\cod{cosine}) y la \textbf{profundidad} de recuperación
(\cod{limit\_per\_entry}), que es cuántos vecinos se traen por query. El cuidado
es que la profundidad son en realidad dos cantidades distintas que el sistema no
debe mezclar: el número de vecinos que se \emph{recuperan}, que cambia qué
candidatos existen y exige volver a ejecutar la recuperación, y el
\emph{corte de evaluación} (\cod{max\_k\_position}), que trunca una lista ya
recuperada y no exige volver a calcular nada. Un guardia específico se niega a
comparar dos resultados cuyas profundidades sean cantidades distintas.}

\term{Generator (etapa 4). Si algún candidato llega sin donante}{Mide si existe
una vía de generación de candidatos que no sea la transferencia por vecindad,
típicamente anotación por dominios de InterPro. Hoy la respuesta es que no
existe, y el esquema lo garantiza en lugar de constatarlo.}

\term{Scoring (etapa 5). Qué ponderación convierte un candidato en una
puntuación}{La fórmula que asigna el número. En el estado actual todos los
números del registro los produjo la vía de respaldo,
$1 - \text{distancia coseno}/2$, que no está registrada como configuración. El
nodo existe para que una ponderación alternativa sea comparable con esa.}

\term{Features (etapa 6). Qué rasgos por candidato entran en un modelo}{El
conjunto de columnas que describe un candidato: identidad de secuencia con el
donante, relación taxonómica, information accretion del término, estadísticos de
la votación de vecinos, y demás. El esquema de rasgos está versionado y agrupa
21 familias sobre 20 conjuntos de columnas independientes. La ablación correcta
retira \emph{columnas}, nunca familias, porque una de las familias (\cod{knn}) es
la unión disjunta de otras dos (\cod{knn\_distance} y \cod{knn\_vote}); las otras
dieciocho son disjuntas dos a dos, y de ahí que 21 familias den 20 conjuntos.}

\term{Reranking (etapa 7). Si un modelo reordena los candidatos}{El paso que
convierte el sistema de un método de recuperación en un método aprendido. Toma
los rasgos de la etapa 6 y reordena. Es el objeto de la segunda pregunta de
investigación del proyecto y hoy no tiene ninguna instancia.}

\term{Combination (etapa 8). Cómo se fusionan dos o más flujos en una respuesta}{
Relevante solo cuando existe más de una vía de generación, es decir, cuando el
nodo generator deja de estar bloqueado.}

\term[def:routing]{Routing (etapa 9). Qué flujo responde a qué panel}{La decisión
de encaminar configuraciones distintas a regiones distintas de la población, en
lugar de elegir una sola configuración para todo. Es lo que el proyecto llama
\emph{canalización por estratificación masiva}, y la sección~\ref{sec:estadistico}
explica por qué es la decisión con más riesgo estadístico de las diez.}

\subsection{Nodo, nivel, arista y firmeza}

\term[def:nodo]{Node}{Una decisión sobre un campo, o sobre un grupo de campos que
no pueden decidirse por separado. Las diez de la sección~\ref{sec:diez} son todas
las que el sistema reconoce, y la lista es fija en el código: una prueba recorre
la especificación, resuelve el constructor de cada nodo por su nombre, comprueba
que hay diez y comprueba que a los diez se les entrega el umbral declarado.}

\term[def:nivel]{Level}{\textbf{Un valor junto con todos los campos que varían
entre las filas con las que se compara.} Esta es la definición más importante del
documento y la fuente del defecto que más veces ha reaparecido en el proyecto.

El motivo es que un nombre corto es casi siempre insuficiente. Si se comparan dos
arms que difieren en el modelo y \emph{también} en la profundidad de
recuperación, entonces \cod{ankh\_large} no nombra un nivel: nombra dos arms
distintos que comparten el nombre del modelo. La comparación parece de un solo
campo y no lo es.

De aquí se sigue una propiedad incómoda pero real: \textbf{la aridad de un nivel
se decide en el momento de la lectura, no en el momento de la declaración}. Sobre
el registro completo hacen falta cinco campos para nombrar un nivel; restringido
a un solo frame seal, \cod{donor\_policy} pasa a ser constante y bastan cuatro.
La consecuencia práctica es que \emph{ninguna cadena de texto almacenada puede
nombrar un nivel}, y por tanto un umbral debe declararse como una tupla de campos
más valores, nunca como un nombre ya compuesto.}

\term{Edge}{Lo que el sistema sabe sobre un nodo: si el artefacto que hace falta
existe, cuántos niveles se han instanciado, cuántos están disponibles, cuántos se
han puntuado, cuántos resultados hay, si el marco fuerza el valor, cuál es el
umbral declarado y si la comparación lo superó.}

\term[def:firmeza]{Strength}{La palabra que dice con qué firmeza se sostiene una
decisión. Hay cinco y se prueban en un orden fijo, que es el argumento del
diseño:

\begin{center}
\begin{tabularx}{\linewidth}{@{}c L{3.2cm} X@{}}
\toprule
\# & palabra & se devuelve cuando\\
\midrule
1 & \cod{blocked}    & el artefacto no existe, o no hay ningún nivel instanciado\\
2 & \cod{chosen} o \cod{inherited} & hay un solo nivel. \cod{chosen} si el marco lo fuerza, \cod{inherited} si nadie lo eligió\\
3 & \cod{unpowered}  & hay dos niveles o más, pero menos de dos puntuados\\
4 & \cod{chosen}     & hay poder, pero no hay umbral declarado, o el umbral no pudo responder\\
5 & \cod{measured}   & hay poder, hay umbral, y la comparación lo superó. Si no lo supera, \cod{chosen}\\
\bottomrule
\end{tabularx}
\end{center}

Dos propiedades del orden importan aguas abajo. \textbf{El poder se prueba antes
que la evidencia}: ninguna diferencia, por grande que sea, puede convencer a un
nodo de que hubo una medida si la comparación no tenía forma de resolver nada.
Y \textbf{\cod{chosen} es un sumidero}: cuatro situaciones distintas caen en esa
palabra, de modo que un lector que solo vea la palabra no puede distinguir una
decisión sin umbral de una que tenía umbral y no lo superó. Publicar el umbral y
el resultado de la comparación junto al nodo es un cambio pequeño y resuelve eso.}

\term[def:suelo]{Floor}{El nivel incumbente contra el que se compara, declarado
\emph{antes} de medir. Es lo único que hace falta desde fuera de las tablas de
resultados para que un nodo pase de \cod{chosen} a \cod{measured}, y por eso es
el punto donde una campaña se vuelve falsable o deja de serlo.}

\term[def:haz]{Survivor set, o haz}{El conjunto de niveles que ninguna celda con
poder llegó a batir. Es el resultado correcto cuando un eje no sella: no un
ganador, ni un empate, sino los que siguen en pie. Un eje de esta campaña ya
produjo exactamente eso, un haz de tres, y la instrumentación actual no sabe
transportarlo porque solo puede leer un umbral de un solo nombre.}

\subsection{La unidad de reporte}
\label{sec:unidad}

\term[def:category]{Category}{La partición por conocimiento previo de la query,
en el momento de la cota inferior de la ventana. Son tres:
\cod{NK} (\emph{no knowledge}), proteínas sin ninguna anotación previa en el
aspecto; \cod{LK} (\emph{limited knowledge}), con anotación en otro aspecto pero
no en este; y \cod{PK} (\emph{prior knowledge}), con anotación previa en este
mismo aspecto. Es la partición que usa CAFA y separa tres problemas distintos:
predecir desde cero, transferir entre aspectos y completar.}

\term[def:aspect]{Aspect}{La sub-ontología: \cod{BPO} (proceso biológico),
\cod{MFO} (función molecular) y \cod{CCO} (componente celular). Son tres grafos
separados con estadística distinta, y promediar entre ellos mezcla tres
problemas.}

\term[def:panel]{Panel}{El producto de category por aspect: nueve celdas.
\textbf{Es la partición obligatoria y nunca se agrega.} Ni la campaña, ni la
superficie que la publica, ni la operación que compara dos arms suman o promedian
los nueve paneles. La operación de comparación se niega explícitamente y nombra
su propio párrafo al negarse.

La razón no es de estilo. Un promedio sin pesos entre poblaciones que difieren en
más de un orden de magnitud es en sí mismo un repesado, y va en la dirección
favorable, porque los estratos más pequeños son los más fáciles. Los nueve
paneles se emiten siempre, con resultado o sin él, para que una corrida parcial no
pueda leerse como completa.

La consecuencia para el reporte es que \textbf{un veredicto es un vector de
nueve componentes y nunca un número}. Cada celda lleva su clave de panel, su
población contada sobre la verdad de la propia ventana y nunca inferida, los
niveles puntuados en orden, y la firmeza del veredicto en esa celda.}

\term[def:estrato]{Stratum}{Un punto en siete ejes, y la unidad por la que se
reporta cuando se quiere más resolución que el panel. Los siete son
\cod{category}, \cod{aspect}, \cod{length}, \cod{homology},
\cod{donor\_evidence}, \cod{taxonomy} y \cod{propagation}.

Los siete no son intercambiables, y la división que importa está escrita en el
código: \cod{category}, \cod{aspect} y \cod{length} son propiedades \emph{de la
query} y se pueden leer siempre; los otros cuatro son propiedades \emph{de una
recuperación} y necesitan un donante efectivamente alineado antes de poder
leerse. Pedir un eje de vecindario sin alineamientos calculados descarta
silenciosamente a la mayoría de las queries y reporta a los supervivientes como
si fuesen la población, que es un error que ya ocurrió.}

\term{Length band}{Banda de longitud de secuencia:
\cod{<=512}, \cod{512-1024}, \cod{1024-2048} y \cod{>2048}. Los cortes son los
límites de contexto de los modelos, no cuantiles del corpus. Una proteína más
larga que el contexto se trunca antes de codificarse, de modo que las bandas
separan secuencias que la representación vio enteras de secuencias que vio en
parte, que es una pregunta distinta de si una proteína es larga. Bandas por
cuantiles mezclarían las dos y se moverían cada vez que se moviese el corpus.}

\term[def:homologia]{Homology band}{Lo cerca que está el donante anotado más
próximo, por identidad de secuencia:
\cod{none}, \cod{<=30}, \cod{30-60}, \cod{60-90} y \cod{>90}.
Los cortes marcan donde la transferencia por similitud cambia de carácter: por
debajo del 30\,\% está la zona de penumbra, donde el alineamiento por sí solo
deja de implicar función compartida; por encima del 90\,\% la transferencia es
casi una copia y dice poco sobre si un método generaliza. \cod{none} es una banda
y no un valor ausente: una query sin ningún donante es una población real y
reportable, y plegarla en la banda más baja escondería los fallos de recuperación
dentro de los resultados de modelado.}

\term{Donor evidence}{Si la anotación del donante más cercano se apoya en
experimento: \cod{exp}, \cod{other} o \cod{none}. Se cruza con la banda de
homología porque las dos juntas son lo que un umbral de homología significa de
verdad. Un donante cercano cuya propia anotación fue inferida por cómputo
transfiere una conjetura, y reportarlo junto a un donante cercano con soporte
experimental promediaría dos afirmaciones distintas en un número.}

\term{Taxonomy band}{A qué distancia está el donante más cercano en el árbol de la
vida: \cod{same}, \cod{close}, \cod{intermediate}, \cod{distant},
\cod{root-only} o \cod{none}. Se resuelve sobre donantes de secuencia distinta a
la query, que es la única lectura bajo la cual \emph{misma especie} es un hecho
sobre el vecindario y no un hecho sobre el protocolo.}

\term[def:propagacion]{Propagation band}{Cuánto tuvo que viajar una anotación para
llegar a la query: \cod{none}, \cod{0}, \cod{0-0.05}, \cod{0.05-0.15} y
\cod{>0.15}. No es la banda de homología. Esa pregunta cuán cerca está el donante
más cercano de cualquier tipo; esta pregunta cuánto más lejos hay que ir para
llegar a un donante cuya propia anotación se apoya en experimento, es decir, la
diferencia entre dos distancias. Una proteína puede estar en un vecindario denso
y aun así lejos de cualquier evidencia experimental, y ese hueco es la distancia
que la anotación recorrió.}

\term[def:floorpop]{Population floor}{El tamaño mínimo que una celda debe
tener para que se le permita votar, derivado de la desviación típica emparejada y
del efecto que se quiere poder detectar. Para el encaminamiento vale
$\lceil (2.8016\,\sigma/0.02)^2 \rceil$. Con un $\sigma$ global de 0.13 da 332,
que es la cifra que la campaña ha usado; con el $\sigma$ implícito de cada panel va
de 120 a 1{,}376, y bajo ese criterio tres de los nueve paneles no alcanzan su
propio suelo. Una celda por debajo del suelo no se reporta como un veredicto débil:
no se reporta.}

\subsection{Los estadísticos}

\term[def:fmax]{\cod{fmax\_w}}{La medida centrada en proteína que usa CAFA. Se
barre un umbral $\tau$, se calculan precisión y exhaustividad promediadas sobre
proteínas y ponderadas por information accretion, y se toma el máximo de la
$F_1$ sobre la curva. \textbf{No admite un error típico de la forma
$\sigma/\sqrt{n}$}, porque es un máximo sobre $\tau$ de una combinación no lineal
de dos medias poblacionales. Cualquier intervalo que se le ponga por esa vía es
inventado.}

\term{\cod{f\_micro\_w}}{La variante agrupada, que acumula los conteos sobre todos
los pares antes de calcular la $F_1$. Da más peso a las proteínas con muchas
anotaciones.}

\term[def:bootstrap]{Bootstrap pareado por proteína}{El estadístico que decide en
esta campaña. Para cada proteína se calcula su propia mejor $F_1$ bajo cada uno de
los dos arms que se comparan, se toma la diferencia por proteína, y se remuestrean
las proteínas con reemplazo para obtener un intervalo de confianza de la media de
esa diferencia. Es \emph{pareado} porque las dos configuraciones se evalúan sobre
las mismas proteínas, lo que elimina la varianza entre proteínas, que es la
dominante.

Este estadístico y \cod{fmax\_w} no son dos formas de medir lo mismo. Premian
comportamientos distintos: \cod{fmax\_w} maximiza una curva agregada y recompensa
predecir poco y con confianza, mientras que el pareado da a cada proteína su
propio óptimo y recompensa cubrir a cada proteína. En los dos ejes que la campaña
ya cerró, los dos estadísticos eligieron ganadores distintos, y en un caso
extremos opuestos de la misma rejilla. De ahí la propuesta de que el estadístico
sea un campo del sello (sección~\ref{sec:marco}).}

\term[def:mde]{MDE, efecto mínimo detectable}{La diferencia más pequeña que una
celda puede distinguir de cero con su población y su varianza. Se lee de la propia
distribución bootstrap y es la cifra que decide si una celda tiene voto. Sobre los
nueve paneles de esta campaña la mediana va de 0.0019 en \cod{PK:BPO} a 0.0256 en
\cod{LK:CCO}, un factor de 13.3; y dentro de un mismo panel varía por un factor de
más de veinte según la comparación, que es la razón por la que se publica como
rango y no como constante.}

\term{Cuantización}{Todo lo almacenado está redondeado a $10^{-4}$. Una diferencia
que se lee como cero desde la base de datos es una \emph{cota superior} del
acuerdo, nunca un cero medido. Junto con el hecho de que el MDE de la celda peor
resuelta es 0.0234, esto significa que una regla basada en estimaciones puntuales
puede decidir a partir de diferencias doscientas veces más pequeñas que lo que el
dato resuelve.}

\section{El grafo de decisiones}
\label{sec:grafo}

\subsection{Forma}

La campaña se representa como un grafo de diez nodos en diez etapas ordenadas,
cada uno con su pregunta y su palabra de firmeza. No es un diagrama ilustrativo:
es un objeto que se calcula a partir del registro, y cada nodo publica junto a su
palabra el valor en el que está parado, campo por campo. Una palabra dice con qué
firmeza se sostiene una decisión y no dice nada sobre \emph{qué} se decidió; un
lector que no vea el valor no puede distinguir un valor por omisión heredado de
una elección deliberada que resulta no estar medida. Las dos leen
\cod{inherited} y solo una de las dos es una sorpresa.

\begin{center}
\begin{tabularx}{\linewidth}{@{}r L{2.6cm} X@{}}
\toprule
etapa & nodo & pregunta\\
\midrule
0 & \cod{frame}       & En qué ventana, pivote y régimen de accretion se lee cada número\\
1 & \cod{substrate}   & En qué representación se calcula el vecindario\\
2 & \cod{bank}        & De qué corpus vienen los donantes, y bajo qué política\\
3 & \cod{retriever}   & Cómo se extraen candidatos del banco, y a qué profundidad\\
4 & \cod{generator}   & Si algún candidato llega sin donante\\
5 & \cod{scoring}     & Qué ponderación convierte un candidato en una puntuación\\
6 & \cod{features}    & Qué rasgos por candidato entran en un modelo\\
7 & \cod{reranking}   & Si un modelo reordena los candidatos\\
8 & \cod{combination} & Cómo se fusionan dos o más flujos en una respuesta\\
9 & \cod{routing}     & Qué flujo responde a qué panel\\
\bottomrule
\end{tabularx}
\captionof{table}{Los diez nodos. El orden es el del flujo de datos.}
\end{center}

\subsection{Cómo se cierra un nodo}

Las transiciones entre palabras son cuatro, y solo la última necesita algo que no
esté en las tablas de resultados:

\begin{center}
\begin{tabularx}{\linewidth}{@{}L{4.6cm} X@{}}
\toprule
transición & lo que hace falta\\
\midrule
\cod{blocked} $\rightarrow$ cualquiera & producir el artefacto que falta. Cada nodo bloqueado publica su propia precondición como un conteo leído del registro\\
\cod{inherited} $\rightarrow$ \cod{unpowered} & instanciar un segundo nivel. Se cuenta de las filas, nunca de una declaración\\
\cod{unpowered} $\rightarrow$ \cod{chosen} & puntuar el segundo nivel\\
\cod{chosen} $\rightarrow$ \cod{measured} & declarar un umbral que nombre un nivel real, y superarlo\\
\bottomrule
\end{tabularx}
\end{center}

\subsection{Por qué el grafo es la estructura y no las fases}

Una organización alternativa, y la que el proyecto usó antes, es por fases
sucesivas: fase A, fase B, y así. Tiene dos defectos frente al grafo.

El primero es que una fase no tiene tipo. Decir que la fase A está terminada no
dice si terminó porque se midió algo o porque se dejó de mirar. Un nodo con una
palabra de firmeza sí lo dice, y además no puede decirlo de más: la palabra se
calcula del registro y no se escribe a mano.

El segundo es que las fases sugieren un orden total donde solo hay un orden
parcial. El nodo \cod{routing} depende de que \cod{substrate} y \cod{retriever}
estén resueltos, pero no depende en absoluto de \cod{generator}. Una fase que
agrupe los tres obliga a esperar por el que no hace falta. El grafo permite que
seis nodos estén \cod{blocked} sin que eso impida avanzar en los otros cuatro,
que es exactamente la situación actual y no es un problema.

\section{El marco: qué hace comparables dos números}
\label{sec:marco}

\subsection{El sello}

Dos números son comparables si y solo si se leyeron bajo las mismas condiciones.
Esas condiciones se resumen en un digest criptográfico que se almacena en cada
resultado, y el material que lo produjo se guarda al lado del digest, de modo que
un sello se puede reconstruir y auditar en lugar de solo verificar.

Hoy el sello cubre seis campos. El diseño propone ocho, y la diferencia está
razonada campo por campo:

\begin{center}
\begin{tabularx}{\linewidth}{@{}L{5.1cm} c X@{}}
\toprule
campo & hoy & razón\\
\midrule
\cod{evaluation\_set\_id}            & sí & la ventana. Sin ella no hay nada que comparar\\
\cod{pivot\_snapshot\_id}            & sí & el grafo bajo el que se resuelve la verdad\\
\cod{information\_accretion\_set\_id}& sí & los pesos por término\\
\cod{max\_terms}                     & sí & cuántos candidatos se conservan por query\\
\cod{max\_distance}                  & sí & el corte de distancia\\
\cod{temporal\_window}               & sí & \textbf{propuesto retirar.} Es texto libre de un campo del arnés, sin restricción de vocabulario. Dos sellos distintos difieren hoy solo en que uno lo tiene nulo y el otro lleva una cadena, y son por lo demás idénticos\\
\cod{th\_step}                       & no & \textbf{propuesto añadir.} Es la resolución de la rejilla de umbrales. Refinarla de 0.01 a 0.0005 invirtió el veredicto en una parte sustancial de las comparaciones de un eje. Los cuatro sellos actuales contienen a la vez la rejilla gruesa y la fina\\
\cod{deciding\_statistic}            & no & \textbf{propuesto añadir.} En los dos ejes ya cerrados el ganador dependió de cuál de los dos estadísticos se leyera, y en un caso los dos eligieron extremos opuestos de la misma rejilla. Si dos números no son comparables cuando difieren en el pivote, tampoco lo son cuando difieren en el estadístico que los decide\\
\bottomrule
\end{tabularx}
\end{center}

\begin{aviso}{Decisión requerida: re-sellar cambia todos los digests}
Añadir o retirar un campo del sello cambia el digest de todos los resultados
existentes. Cualquier digest citado en un documento o en una nota queda obsoleto,
y las filas antiguas del catálogo de sellos sobreviven junto a las nuevas porque
la inserción ignora conflictos. Es un cambio sobre un objeto declarado y necesita
aprobación explícita, no silencio.
\end{aviso}

\subsection{La ventana de selección y el holdout}

La campaña selecciona sobre $220 \rightarrow 227$ y no lee nada por encima de 227
del lado nuevo. La regla operativa es más amplia que un intervalo fijo y se
enuncia así: \textbf{nada en la release 227 o posterior, del lado nuevo, informa
ninguna elección}.

Conviene separar dos cosas que se confunden con facilidad, porque de la distinción
depende que el diseño sea más fuerte de lo que parece:

\begin{itemize}[leftmargin=1.4em,itemsep=0.3em]
  \item \textbf{Puntuarse a sí mismo sobre datos posteriores} para elegir una
        configuración es fuga, y está prohibido. Un solo punto de validación
        interna se designa antes de empezar, y designarlo no autoriza a leerlo
        hasta el final.
  \item \textbf{Enviar predicciones a una evaluación externa} no es fuga. En CAFA
        y en LAFA se predice sobre proteínas cuya anotación \emph{todavía no
        existe}, y los organizadores puntúan después contra anotaciones que se
        acumulan tras el cierre de envío. La fuga es imposible por construcción,
        porque no hay nada que mirar.
\end{itemize}

De ahí una propiedad del diseño que merece decirse en voz alta: \textbf{si la
validación final es un despliegue externo, el problema del holdout se resuelve
solo}. El holdout verdadero es el futuro, y nadie puede espiarlo. Toda la
disciplina sobre 227 se aplica únicamente a los números \emph{internos} que la
campaña cita para justificar sus propias elecciones.

La consecuencia práctica es la inversa de lo que suele suponerse: el rigor interno
no es menos importante porque haya validación externa, es \emph{más} importante.
El número de envíos externos es limitado, y cada configuración enviada que fue
elegida sobre ruido es un envío desperdiciado. El trabajo del intervalo de
confianza es precisamente garantizar que lo que se despliega es lo que de verdad
ganó.

\subsection{La protección por ausencia}

El mecanismo más fuerte que el proyecto tiene contra la fuga no es un guardia, es
una ausencia: las releases que no deben mirarse no están cargadas en la base. Si
el dato no está, no se puede mirar por accidente, ni por una consulta distraída,
ni por un trabajo mal parametrizado. Reponer la release del holdout exige una
descarga de 14.4~GB, y esa fricción es lo que convierte el error potencial en un
acto deliberado.

\section{La superficie de parámetros}
\label{sec:parametros}

Esta sección enumera \emph{todos} los parámetros que existen, con el nodo al que
pertenecen, dónde viven, su tipo y los niveles que se han instanciado de hecho. La
enumeración se obtuvo de la base de datos y no del código, porque lo que decide
qué se ha explorado son las filas y no las firmas de las funciones.

Un parámetro que aparece con un solo nivel instanciado no es un parámetro
explorado: es una constante que nadie eligió, y esa es exactamente la lectura que
la palabra \cod{inherited} le da.

\subsection{Predicción: quince parámetros}

\begin{center}
\small
\begin{longtable}{@{}L{4.9cm} L{1.6cm} L{1.6cm} L{5.6cm}@{}}
\toprule
parámetro & nodo & tipo & niveles instanciados\\
\midrule
\endfirsthead
\toprule
parámetro & nodo & tipo & niveles instanciados\\
\midrule
\endhead
\cod{embedding\_config\_id}      & substrate & uuid    & \textbf{25}. La rejilla completa del sustrato, desglosada en la tabla~\ref{tab:sustrato}\\
\cod{metric}                     & retriever & enum    & 1: \cod{cosine}\\
\cod{limit\_per\_entry}          & retriever & entero  & 1: 200, en los 86 prediction sets. El valor 30 aparece en un único trabajo que falló y no produjo ninguno. Es la profundidad de recuperación, no el corte de evaluación\\
\cod{search\_backend}            & retriever & enum    & 1: \cod{numpy}. El KNN corre siempre en CPU; el valor por omisión \cod{auto} elegía GPU en silencio, que es el único camino que falla\\
\cod{annotation\_set\_id}        & bank      & uuid    & 1: GOA 220. Lo fija el frame\\
\cod{donor\_policy}              & bank      & objeto  & 2: la restrictiva y la vacía. La restrictiva admite once códigos experimentales y de alto rendimiento más \cod{IC} y \cod{TAS}, trece en total; la vacía no filtra por código de evidencia y es por tanto \textbf{la más permisiva de las dos}. El filtro solo se aplica \cod{if codes}, de modo que la lista restringe\\
\cod{exclude\_self\_neighbour}   & retriever & booleano& 2. Retira el auto-acierto a distancia cero\\
\cod{expand\_votes\_to\_ancestors}& retriever& booleano& 2. Propaga el voto de un donante a los ancestros del término\\
\cod{aspect\_separated\_knn}     & retriever & booleano& 2. Recupera vecinos por aspecto en lugar de una sola vez\\
\cod{ontology\_snapshot\_id}     & frame     & uuid    & 1. El grafo bajo el que se resuelven los términos del banco\\
\cod{query\_set\_id}             & frame     & uuid    & 1. El conjunto de queries de la campaña\\
\cod{compute\_alignments}        & features  & booleano& 1: \cod{false} en los 86 prediction sets. El nivel \cod{true} solo aparece en trabajos que no produjeron ninguna fila; la identidad disponible viene del backfill de \cod{sequence\_alignment}\\
\cod{compute\_taxonomy}          & features  & booleano& 1: \cod{false} en los 86 prediction sets, igual que la anterior\\
\cod{compute\_reranker\_features}& features  & booleano& 1: \cod{false} en los 86 prediction sets, igual que las dos anteriores\\
\cod{batch\_size}                & ninguno   & entero  & 1: 1024. Parámetro de ejecución, no de método. No debe entrar en la identidad de un nivel\\
\bottomrule
\caption{Los quince parámetros de \cod{predict\_go\_terms}, con sus niveles
instanciados sobre 69 trabajos.}
\end{longtable}
\end{center}

\begin{aviso}{Las banderas \cod{compute\_} son precondiciones, no ablaciones}
Las tres banderas \cod{compute\_} no cambian la predicción: cambian qué se
\emph{registra} sobre ella. Apagarlas no es una ablación del método, es renunciar
a poder estratificar después. \cod{compute\_alignments} apagado hace ilegibles
cuatro de los siete ejes de estrato, porque los ejes de vecindario necesitan un
donante alineado. La ablación de rasgos, que sí es una decisión de método,
pertenece al nodo \cod{features} y se hace retirando columnas del esquema de
rasgos, nunca apagando el cálculo.
\end{aviso}

\subsection{Evaluación: siete parámetros}

\begin{center}
\small
\begin{tabularx}{\linewidth}{@{}L{5.2cm} L{1.6cm} X@{}}
\toprule
parámetro & nodo & niveles instanciados\\
\midrule
\cod{prediction\_set\_id}      & --       & 46. El arm que se evalúa\\
\cod{evaluation\_set\_id}      & frame    & 2. Las dos variantes de la ventana $220 \rightarrow 227$\\
\cod{information\_accretion\_set\_id} & frame & 1\\
\cod{max\_k\_position}         & retriever & \textbf{7}: 1, 2, 3, 5, 10, 20, 30. Es el corte de evaluación, que trunca una lista ya recuperada\\
\cod{th\_step}                 & frame    & 2: 0.0005 y 0.002, más la ausencia del campo, que toma el valor por omisión 0.01\\
\cod{temporal\_window}         & frame    & 1: la cadena \cod{SELECT\_220\_227}. Etiqueta del arnés\\
\cod{frame}                    & --       & 1: la cadena \cod{internal}. Etiqueta del arnés\\
\bottomrule
\end{tabularx}
\end{center}

\subsection{El sustrato, desglosado}

De los once campos que determinan un substrate, solo siete varían de hecho en el
registro, y de esos siete cuatro solo distinguen la familia de modelos
preentrenados de las tres codificaciones aprendidas. El eje real lo llevan
\cod{model\_name} (12 valores), \cod{model\_backend} (7) y \cod{layer\_indices} (12). Es decir: \textbf{el eje del sustrato, tal
como está instanciado, es modelo por capa}.

\begin{center}
\small
\begin{tabularx}{\linewidth}{@{}L{5.3cm} L{2.8cm} L{2.4cm} X@{}}
\toprule
\cod{model\_name} & \cod{backend} & capas & notas\\
\midrule
\cod{facebook/esm2\_t6\_8M\_UR50D}   & \cod{esm}    & 0 & el más pequeño, usado como suelo\\
\cod{facebook/esm2\_t33\_650M\_UR50D}& \cod{esm}    & 0, 11, 22, 33 & rejilla de capas\\
\cod{facebook/esm2\_t36\_3B\_UR50D}  & \cod{esm}    & 0 & \\
\cod{esmc\_600m}                     & \cod{esm3c}  & 0, 12, 24, 35 & rejilla de capas\\
\cod{ElnaggarLab/ankh-base}          & \cod{ankh}   & 0, 10, 16, 32, 48 & las capas 10, 16 y 32 \textbf{no son evaluables}: puntuación constante\\
\cod{ElnaggarLab/ankh-large}         & \cod{ankh}   & 0 & \\
\cod{Rostlab/prot\_t5\_xl\_half}     & \cod{t5}     & 0, 8, 16, 24 & rejilla de capas\\
\cod{Rostlab/ProstT5}                & \cod{t5}     & 0 & \\
\cod{mila-intel/ProtST-esm1b}        & \cod{protst} & 0 & \\
\cod{rung2-dense}, \cod{rung2-pooled}& \cod{learned-code} & --- & codificaciones aprendidas. \cod{max\_length} 1022, sin normalizar\\
\cod{rung2-residue}                  & \cod{residue-sparse} & --- & idem\\
\bottomrule
\end{tabularx}
\captionof{table}{Las 25 configuraciones de sustrato registradas, en siete
familias de modelo y doce nombres de modelo distintos, nueve preentrenados y tres
aprendidos. \cod{pooling} es \cod{mean} en las veintidós preentrenadas y
\cod{learned} o \cod{residue-sparse-mean} en las tres aprendidas;
\cod{max\_length} es 2048 en las preentrenadas y 1022 en las aprendidas.}
\label{tab:sustrato}
\end{center}

\begin{aviso}{Un confundido que invalida el eje del sustrato tal como está hoy}
Doce de las configuraciones tienen 528{,}294 embeddings almacenados y once tienen
60{,}797, un factor de \textbf{8.69}. Una comparación entre un arm con el banco
grande y otro con el banco pequeño mide dos cosas a la vez, la representación y el
tamaño del banco de donantes, y ninguna cantidad posterior las separa. El eje del
sustrato no puede producir un veredicto hasta que los bancos se igualen. El coste
de igualarlos está en la sección~\ref{sec:presupuesto} y es asumible.
\end{aviso}

\section{La unidad de reporte}
\label{sec:unidad-reporte}

\subsection{Dos niveles, no uno}

El reporte tiene dos niveles y conviene no mezclarlos, porque tienen funciones
distintas.

El \textbf{panel} es la partición obligatoria: category por aspect, nueve celdas,
siempre las nueve, nunca agregadas. Su función es impedir que un promedio esconda
que el método funciona en un problema y no en otro. Las nueve celdas son nueve
problemas distintos, no nueve muestras del mismo.

El \textbf{estrato} es el refinamiento dentro del panel: los siete ejes de
D\ref{def:estrato}. Su función es distinta, y es la que interesa al proyecto:
describir a cada proteína \emph{por su posición en el vecindario}, no solo por su
conocimiento previo. Una proteína corta con un donante casi idéntico y una
proteína larga sin ningún donante anotado son dos problemas de recuperación
completamente distintos que el panel por sí solo no separa.

El mecanismo que lo hace posible sin violar la regla de no agregación es que el
estrato \textbf{viaja dentro de la clave del panel}. Una celda se identifica como
\cod{NK:MFO@length=512-1024}, de modo que una restricción por estrato produce
nueve celdas restringidas y nunca un número.

\begin{aviso}{La pertenencia al estrato se toma del brazo base, no del comparado}
Si cada arm decidiera la banda de sus propias queries, dos arms estarían
comparándose sobre poblaciones distintas y la diferencia medida incluiría el
cambio de población. La pertenencia se resuelve una vez, sobre el brazo base, y se
aplica a los dos lados.
\end{aviso}

\subsection{El límite real: granularidad contra población}

Aquí está la tensión que el diseño tiene que resolver, y es aritmética.

El cruce completo de los siete ejes tiene
$3 \times 3 \times 4 \times 5 \times 3 \times 6 \times 5 = 16{,}200$ celdas. La
campaña dispone de unas 30{,}427 unidades emparejadas en total. Es decir, el cruce
completo promedia menos de dos proteínas por celda, y el suelo de población para
que una celda tenga voto es 332 con el $\sigma$ global, y entre 120 y 1{,}376 con el
de cada panel (D\ref{def:floorpop}). \textbf{El cruce de
siete ejes está vacío en casi todas partes}, y no por falta de cómputo: por falta
de proteínas.

Medido sobre la campaña actual, cruzando el panel con una sola banda de longitud,
sobreviven \textbf{17 celdas} bajo el suelo global de 332, y solo \textbf{11} bajo
el suelo por panel. Y las once viven en cinco paneles: cuatro en \cod{PK:BPO},
tres en \cod{PK:MFO}, dos en \cod{PK:CCO}, una en \cod{LK:BPO} y una en
\cod{NK:BPO}. \textbf{Cuatro paneles no aportan ninguna región}: \cod{NK:MFO},
\cod{NK:CCO}, \cod{LK:MFO} y \cod{LK:CCO}. Dicho de otro modo, encaminar por
longitud sobre un suelo defendible solo puede encaminar dentro de \cod{PK} y de
\cod{BPO}: no puede encaminar \cod{MFO} ni \cod{CCO} fuera de \cod{PK}.

De esto se siguen tres reglas de diseño:

\begin{enumerate}[leftmargin=1.5em,itemsep=0.35em]
  \item \textbf{El panel es obligatorio; el refinamiento es de un eje a la vez.}
        Panel por longitud, o panel por homología, pero no panel por longitud por
        homología, que no tiene población.
  \item \textbf{El eje del refinamiento se declara antes de la corrida.} Si se
        prueban varios y se reporta el que separó, eso es seleccionar sobre ruido
        con otro nombre. La sección~\ref{sec:estadistico} cuantifica cuánto ruido.
  \item \textbf{El cruce completo es una herramienta de descripción del campeón,
        no de selección.} Una vez elegida una configuración, describirla sobre los
        siete ejes es informativo y no decide nada. Usarla para elegir es lo que no
        se puede hacer.
\end{enumerate}

\subsection{Los dos ejes que el proyecto quiere, y su estado}

\begin{center}
\begin{tabularx}{\linewidth}{@{}L{2.3cm} L{3.4cm} X@{}}
\toprule
eje & precondición & estado\\
\midrule
\cod{length}    & ninguna. Es una propiedad de la query & disponible sobre todas las queries. Cuatro bandas por los límites de contexto\\
\cod{homology}  & \cod{compute\_alignments} activo y el backfill de identidad completo & disponible. Medido tres veces por panel y solo bajo dos de los sellos, lo que es poco\\
\bottomrule
\end{tabularx}
\end{center}

La banda de homología es la que más información aporta sobre el método, porque es
la que separa el régimen donde la transferencia por similitud es casi una copia del
régimen donde deja de implicar función compartida. Un resultado ya medido en la
campaña lo ilustra: la política sin filtro, que es la más permisiva de las dos,
gana en \textbf{196 de las 197 lecturas que resuelven} (la excepción es
\cod{LK:MFO} en la banda \cod{>90}), y gana con un efecto entre un 69 y un 78\,\%
mayor por debajo del 30\,\% de identidad que por encima del 90\,\%, según cómo se
agregue. La banda \cod{none} \textbf{no está medida}: sus seis comparaciones
fallaron, de modo que la homología está medida sobre cuatro bandas y no cinco. El mismo veredicto, con una lectura mucho más informativa.

\subsection{Población y poder por panel}

Toda publicación de un veredicto lleva al lado esta tabla, porque el poder de
resolución de los nueve paneles abarca casi un orden de magnitud y una celda que
no puede resolver nada no debe leerse como una celda que no encontró nada.

\begin{center}
\small
\begin{tabularx}{\linewidth}{@{}l R{1.7cm} R{1.9cm} R{2.5cm} R{1.5cm} R{1.5cm} X@{}}
\toprule
panel & $n$ & semiamp. IC & MDE, rango & $\sigma$ impl. & suelo & vota\\
\midrule
\cod{PK:BPO} & 13{,}876 & 0.00130 & 0.0001 - 0.0041 & 0.0782 &   120 & sí\\
\cod{PK:MFO} &  4{,}947 & 0.00421 & 0.0005 - 0.0121 & 0.1511 &   448 & sí\\
\cod{PK:CCO} &  4{,}872 & 0.00466 & 0.0004 - 0.0127 & 0.1659 &   541 & sí\\
\cod{NK:BPO} &  1{,}509 & 0.00951 & 0.0010 - 0.0222 & 0.1884 &   697 & sí\\
\cod{LK:BPO} &  1{,}214 & 0.01102 & 0.0011 - 0.0292 & 0.1958 &   753 & sí\\
\cod{NK:MFO} &  1{,}129 & 0.01544 & 0.0018 - 0.0339 & 0.2647 & 1{,}376 & \textbf{no}\\
\cod{NK:CCO} &  1{,}116 & 0.01520 & 0.0014 - 0.0413 & 0.2591 & 1{,}317 & \textbf{no}\\
\cod{LK:MFO} &     943  & 0.01375 & 0.0016 - 0.0348 & 0.2154 &   911 & sí\\
\cod{LK:CCO} &     821  & 0.01735 & 0.0016 - 0.0431 & 0.2536 & 1{,}263 & \textbf{no}\\
\bottomrule
\end{tabularx}
\captionof{table}{Poder de resolución por panel, sobre la variante
\cod{b7cfed9a}. Todas las columnas salen del registro: $n$ y la semiamplitud del
intervalo bootstrap de los eventos \cod{pairing} y \cod{panel}, el rango del MDE
del propio campo \cod{minimum\_detectable\_effect} sobre las lecturas que
resuelven, y $\sigma$ implícita despejada de la semiamplitud. El MDE se da como
rango porque no es una constante del panel: hay uno por comparación.}
\end{center}

\begin{aviso}{Tres de los nueve paneles no alcanzan su propio suelo}
La columna de suelo aplica la misma regla que el suelo global,
$\lceil (2.8016\,\sigma/0.02)^2 \rceil$, pero con el $\sigma$ implícito de cada
panel en vez de un 0.13 único. Bajo ese criterio \cod{NK:MFO}, \cod{NK:CCO} y
\cod{LK:CCO} \textbf{no pueden votar}: su población está por debajo del mínimo que
su propia dispersión exige. Dos de los tres son \cod{NK}, que es la categoría que
el proyecto quiere titular.

No es un defecto de la tabla, es lo que la tabla mide. Un suelo global de 332
dejaba votar a las nueve porque suponía una dispersión que ocho de los nueve
paneles no tienen.
\end{aviso}

La consecuencia que gobierna la regla de cierre está en esta tabla. Las cinco
celdas más pequeñas son \emph{mayoría de las nueve} y sostienen solo el 17.2\,\%
de la evidencia. La razón entre el error típico mayor y el menor es 13.4. Una regla
que sellase por mayoría simple de los nueve paneles podría cerrar un eje, convertirse
en el umbral del siguiente y propagarse aguas abajo con el voto de la sexta parte
peor resuelta del registro.

\section{El estadístico y la regla de cierre}
\label{sec:estadistico}

\subsection{Para qué sirve un intervalo de confianza en esta campaña}

La objeción natural al intervalo de confianza es que añade ruido a la lectura. Es
la objeción correcta y la respuesta es que el intervalo no añade el ruido: lo
\emph{mide}. Sin él, el ruido sigue ahí y no se ve. Pero conviene ser preciso sobre
cuál es su función aquí, porque no es la función habitual en un artículo.

\textbf{El intervalo es un instrumento de selección, no de publicación.} Su
trabajo es decidir cuál de varias configuraciones avanza a la siguiente etapa. No
hace falta publicarlo para competir en una evaluación externa, donde el árbitro es
la tabla de resultados del organizador. Su valor está enteramente del lado
interno, y se concreta en cuatro cosas.

\subsubsection*{Primero: la cuantización hace que una estimación puntual decida
por debajo de lo que el dato resuelve}

Todo lo almacenado está redondeado a $10^{-4}$. El MDE de la celda peor resuelta es
0.0234. La razón es \textbf{234}. Una regla que compare dos máximos y se quede con
el mayor decide, en el caso peor, a partir de diferencias doscientas treinta veces
más pequeñas que lo que la población de esa celda puede distinguir de cero. Eso no
es una comparación estricta: es una moneda al aire con cuatro decimales.

\subsubsection*{Segundo: el número de oportunidades de acertar por casualidad es
grande, y crece con la ambición del diseño}

La canalización (D\ref{def:routing}) es la decisión donde esto muerde. Encaminar
configuraciones distintas a regiones distintas significa que el arm canalizado
\emph{elige una vez por región} mientras el arm base elige una sola vez. Con quince
regiones, el canalizado tiene quince oportunidades de parecer mejor por azar. Con
nueve paneles y varios arms, el número de contrastes llega al millar sin esfuerzo.

Esto no es una objeción a la canalización: es su condición de validez. El arm
canalizado gana con afinidad real \emph{cero} si no se le exige nada más. Lo único
que lo repara es el \textbf{ajuste cruzado}: la política de encaminamiento se ajusta
sobre unos pliegues y se evalúa sobre otros, de modo que la elección por región no
se puede premiar a sí misma. El intervalo es lo que permite comprobar que el
veredicto fuera de pliegue coincide con el veredicto dentro de pliegue; sin él, las
dos cifras son dos números que difieren y no hay forma de decir si la diferencia
importa.

\subsubsection*{Tercero: convierte \emph{esta celda no puede decir nada} en un
resultado escribible}

El vocabulario que la operación de comparación ya emite tiene seis palabras, y dos
de ellas son las que faltan en cualquier regla de estimación puntual:
\cod{null\_with\_power}, que es una celda con población suficiente cuyo intervalo
cubre el cero, y \cod{null\_unread}, que es una celda que no se pudo leer. La
distinción es la diferencia entre \emph{buscamos y no hay efecto} y \emph{no
pudimos buscar}, y una campaña que no la puede escribir tiene que convertir las dos
en un voto, en cualquier dirección, que es peor que abstenerse.

\subsubsection*{Cuarto: los envíos externos son un recurso escaso}

La validación final del proyecto es un despliegue en LAFA, y a futuro una
participación en CAFA. El número de configuraciones que se pueden enviar es
limitado y el ciclo de realimentación es de meses. Cada configuración enviada que
fue elegida sobre ruido es un envío desperdiciado y una espera perdida. La
validación externa no reduce la necesidad de rigor interno: la aumenta, porque
encarece el error.

\subsection{Qué estadístico decide}

Los dos candidatos no son dos formas de medir lo mismo:

\begin{center}
\begin{tabularx}{\linewidth}{@{}L{3.2cm} X X@{}}
\toprule
 & \cod{fmax\_w} & bootstrap pareado\\
\midrule
qué optimiza & una curva agregada sobre toda la población & el óptimo propio de cada proteína\\
qué premia   & predecir poco y con confianza & cubrir a cada proteína\\
admite $\sigma/\sqrt{n}$ & \textbf{no}. Es un máximo sobre $\tau$ de una combinación no lineal de dos medias & sí, por remuestreo de proteínas\\
comparable con CAFA & sí, es la medida del reto & no directamente\\
\bottomrule
\end{tabularx}
\end{center}

La consecuencia medida es que en los dos ejes que la campaña ya cerró los dos
estadísticos eligieron ganadores distintos, y en uno de los dos eligieron extremos
opuestos de la misma rejilla monótona. No hay respuesta independiente del
estadístico, y por tanto \textbf{el estadístico tiene que declararse, y entrar en
el sello}.

La recomendación es usar los dos con papeles separados y distintos:

\begin{itemize}[leftmargin=1.4em,itemsep=0.3em]
  \item \textbf{El bootstrap pareado decide} qué configuración avanza. Es la
        elección interna, y es la que necesita intervalo.
  \item \textbf{\cod{fmax\_w} se reporta} porque es la medida del reto y es lo que
        hace la cifra comparable con la literatura. Se publica como estimación
        puntual, sin intervalo inventado.
\end{itemize}

\subsection{La regla de cierre}

\begin{aviso}{La regla propuesta, y la que hay que sustituir}
La regla en vigor en el código compara el máximo de \cod{f\_micro\_w} del rival
contra el del incumbente, panel por panel, y exige unanimidad sobre los paneles que
testifican. Son estimaciones puntuales de una métrica cuantizada a $10^{-4}$, sin
intervalo. Además conviven tres reglas incompatibles: esa, una mayoría de los nueve
escrita en un plan, y el bootstrap con IC del 95\,\% escrito en las declaraciones.
Hay que quedarse con una.
\end{aviso}

La regla recomendada tiene tres cláusulas:

\begin{enumerate}[leftmargin=1.5em,itemsep=0.35em]
  \item \textbf{Una celda vota solo si su intervalo del 95\,\% contra el incumbente
        excluye el cero.} Una celda cuyo intervalo cubre el cero reporta
        \cod{null\_with\_power} o \cod{null\_unread} y \emph{no vota en ninguna
        dirección}.
  \item \textbf{Un eje sella solo si un nivel gana en todas las celdas que votan.}
  \item \textbf{Si no, el eje transporta el conjunto supervivente}
        (D\ref{def:haz}): los niveles que ninguna celda con voto llegó a batir.
\end{enumerate}

\subsection{El caso adverso, y por qué la regla necesita tres cláusulas más}

La regla anterior, tal como se enunció, tiene un agujero lógico y un modo de fallo
que conviene poner por escrito antes de adoptarla.

\textbf{El agujero.} Si ninguna celda vota, \emph{un nivel gana todas las celdas
que votan} es vacuamente cierto para todos los niveles a la vez. La regla no define
cuórum, de modo que en el régimen de efectos pequeños no sella: se rompe.

\textbf{El caso adverso, cuantificado con la tabla del propio documento.} Un efecto
uniforme de $+0{,}004$ solo es resoluble en \cod{PK:BPO}, cuyo MDE es $0{,}0032$;
por debajo de $0{,}0126$ ninguna celda \cod{NK} ni \cod{LK} puede votar. Es decir,
en ese régimen la regla degenera en \emph{decide \cod{PK:BPO}}, que es la categoría
\textbf{más fácil} y no la que el proyecto quiere titular. Peor: un efecto real y
pequeño con el mismo signo en 9 de 9 paneles y ningún intervalo excluyendo el cero
da cero votos y no sella, cuando una prueba de signos sobre nueve paneles da
$p = 0{,}004$. \textbf{La regla pierde una decisión correcta que la consistencia sí
sostiene}, y es exactamente la forma de la evidencia que la campaña ya produjo
(9 de 9 celdas, 36 de 36 contrastes).

\textbf{Y la multiplicidad corta en los dos sentidos.} El argumento a favor del
intervalo la invoca, y luego no se la aplica a sí mismo: al 95\,\% por celda, la
probabilidad de que al menos una de nueve vote en falso es del 37\,\%, y con
diecisiete regiones del 58\,\%. Bajo unanimidad, \emph{una sola celda falsa sella o
veta ella sola}. El ajuste cruzado arregla la selección, no la tasa de error por
familia.

Las tres cláusulas que faltan:

\begin{enumerate}[leftmargin=1.5em,itemsep=0.3em]
  \item \textbf{Cuórum.} Un eje no sella si votan menos de tres celdas. Con menos,
        el resultado es \emph{no resuelto}, que es distinto de \emph{no separó}.
  \item \textbf{Segunda vía por consistencia de signo.} Si ninguna celda vota pero
        las nueve tienen el mismo signo, el eje sella con una marca más débil, y la
        marca se publica: sellado por consistencia, no por separación.
  \item \textbf{Nivel por celda corregido.} Con nueve celdas, el nivel por celda
        sube al 99{,}4\,\% para mantener el 5\,\% por familia. Es conservador y hay
        que decir que lo es.
\end{enumerate}

Y el coste que el argumento a favor del intervalo no menciona: \textbf{no sellar
multiplica la rejilla del eje siguiente por el tamaño del haz}, hoy tres, y sin
sello no hay campeón que congelar en la etapa final. El argumento de que los envíos
externos son escasos vale igual en la dirección contraria.

\subsection{Resumen de la regla}

Arregla los dos modos de fallo a la vez: una celda sin poder no puede sostener una
mayoría ni puede vetarla. Es el estadístico que las declaraciones ya nombran y que
los 4{,}279 trabajos de comparación ya ejecutaron. Y es la única opción bajo la cual
\emph{el eje no eligió} es un resultado escribible, que es precisamente lo que uno
de los ejes de esta campaña produjo.

El coste hay que decirlo sin suavizarlo: esto vuelve las filas emparejadas por
proteína imprescindibles para cualquier sellado, y hoy viven solo en un registro de
eventos sin tipar y con borrado en cascada. Promoverlas a una tabla propia deja de
ser limpieza y se convierte en precondición de cerrar cualquier nodo.

\section{La misma campaña a dos escalas}
\label{sec:escalas}

Esta sección es el objeto del informe. La campaña que se ejecuta hoy está acotada
por la potencia de una máquina de cómputo; la que el centro podría ejecutar sobre
cinco máquinas es la misma campaña con una rejilla más ancha. Lo que sigue separa
lo que cambia de lo que no.

\subsection{Lo que no cambia con la escala}

Nada de lo definido en las secciones 2 a 7. En concreto:

\begin{itemize}[leftmargin=1.4em,itemsep=0.25em]
  \item \textbf{Los diez nodos} y su orden. Son propiedades del método, no del
        presupuesto.
  \item \textbf{Las definiciones}, incluida la de nivel, que es la que impide que
        una comparación parezca de un campo cuando varían dos.
  \item \textbf{El marco y su sello.} Un sello no es más barato ni más caro.
  \item \textbf{La unidad de reporte.} Nueve paneles, nunca agregados. Más cómputo
        no produce más proteínas, de modo que el suelo de población por celda es
        exactamente el mismo en las dos escalas.
  \item \textbf{La regla de cierre} y el estadístico que decide.
\end{itemize}

El punto sobre la población merece énfasis porque es contraintuitivo:
\textbf{la escala de cómputo no compra resolución estadística}. El MDE de
\cod{LK:CCO} es 0.0234 porque esa celda tiene 821 proteínas emparejadas, y tendrá
821 con una máquina y con cincuenta. Solo una ventana distinta, o más queries,
cambian eso.

\subsection{Lo que sí cambia}

Dos cosas, y solo dos.

\subsubsection*{La anchura de la rejilla}

Con poco cómputo la rejilla se explora en \textbf{cascada}: se decide un nodo, se
fija su ganador o su haz, y se pasa al siguiente con los demás campos quietos. Con
mucho cómputo se puede explorar \textbf{cruzada}: el producto cartesiano de varios
nodos a la vez.

La diferencia importa porque la cascada supone que las decisiones no interactúan, y
esa suposición es falsa en general. Si el mejor sustrato bajo una política de
donantes permisiva no es el mejor bajo una restrictiva, la cascada lo pierde y el
cruce lo encuentra. \textbf{Lo que el centro compraría con cinco máquinas es
precisamente la capacidad de detectar interacciones entre nodos}, que es la única
cosa que la versión acotada no puede ver por construcción.

\begin{center}
\begin{tabularx}{\linewidth}{@{}L{3.4cm} X X@{}}
\toprule
 & cascada (hoy) & cruce (a escala)\\
\midrule
arms de predicción & 25 del sustrato, luego 4 de política: \textbf{29}. Los cortes no añaden arms, truncan una lista ya recuperada & $22 \times 16 = \mathbf{352}$\\
supuesto           & los nodos no interactúan & ninguno\\
lo que se pierde   & interacciones entre nodos & nada de la rejilla\\
lo que se gana     & tres órdenes de magnitud menos de comparaciones & la interacción, si existe\\
\bottomrule
\end{tabularx}
\end{center}

\subsubsection*{El número de niveles que se pueden comparar a la vez}

Y aquí está el resultado que más conviene que el centro lea, porque va en contra de
la intuición de que más máquinas permiten comparar todo con todo.

Comparar todos los pares de 352 arms son $\binom{352}{2} = 61{,}776$ comparaciones.
Al coste medido de 0.6 minutos cada una son \textbf{798 horas}, que es más que el
coste de producir los 352 arms. Pero el problema no es el tiempo: es que 61{,}776
contrastes sobre nueve paneles son más de setecientas mil pruebas, y con esa
multiplicidad el máximo de una colección de diferencias ruidosas es grande aunque
todas las diferencias verdaderas sean cero.

\begin{aviso}{La cascada no es una concesión al presupuesto, es una necesidad estadística}
Con cómputo infinito se seguiría comparando contra un incumbente declarado y no
todos contra todos, porque el problema de la comparación múltiple no lo resuelve el
cómputo. Lo que cambia con la escala es \emph{cuántos arms se producen} y por tanto
cuántos candidatos entran en la cascada, no la forma de la comparación. El diseño
es el mismo a las dos escalas; a escala grande simplemente entra más gente al
torneo.
\end{aviso}

\subsection{Resumen del reparto}

\begin{center}
\begin{tabularx}{\linewidth}{@{}X c c@{}}
\toprule
 & acotada & a escala\\
\midrule
Nodos de decisión                & 10 & 10\\
Definiciones y marco             & idénticos & idénticos\\
Paneles                          & 9, sin agregar & 9, sin agregar\\
Suelo de población por celda     & 332 & 332\\
Estadístico y regla de cierre    & idénticos & idénticos\\
\midrule
Arms de predicción               & 29 & 352\\
Exploración                      & cascada & cruce\\
Interacciones entre nodos        & no observables & observables\\
Forma de la comparación          & contra incumbente & contra incumbente\\
Bancos de donantes igualados     & pendiente & sí\\
Escalera de recencia del banco   & 2 de 13 releases & 13 de 13\\
Encaminamiento con ajuste cruzado& pendiente & sí\\
\bottomrule
\end{tabularx}
\end{center}

\section{Presupuesto de cómputo}
\label{sec:presupuesto}

\subsection{Costes unitarios medidos}

Todas las cifras salen del registro de trabajos de la campaña actual, sobre
trabajos terminados con éxito. La mediana es la cifra útil; el máximo está para que
nadie planifique con la mediana sola.

\begin{center}
\small
\begin{tabularx}{\linewidth}{@{}L{5.0cm} r r r r@{}}
\toprule
operación & $n$ & mediana & máximo & total\\
\midrule
\cod{predict\_go\_terms}       &    51 & 68.3 min & 377.8 min & 79.6 h\\
\cod{compare\_paired\_panels}  & 4{,}279 &  0.6 min &   8.3 min & 52.1 h\\
\cod{run\_cafa\_evaluation}    &   315 &  2.9 min &  48.1 min & 29.2 h\\
\cod{compute\_embeddings}      &    12 & 84.3 min & 154.6 min & 18.3 h\\
\cod{stratify\_evaluation}     &   309 &  2.9 min &  15.2 min & 18.2 h\\
\cod{load\_goa\_annotations}   &     2 & 128.8 min & 165.6 min &  4.3 h\\
\midrule
\multicolumn{4}{@{}l}{\textbf{total medido de la campaña hasta hoy}} & \textbf{202 h}\\
\bottomrule
\end{tabularx}
\captionof{table}{Coste unitario por operación. \cod{compute\_embeddings} es el
único paso que usa GPU; el KNN corre siempre en CPU.}
\end{center}

Dos observaciones sobre estas cifras antes de proyectarlas. La primera es que
\cod{predict\_go\_terms} tiene una cola larga, de 6.7 a 377.8 minutos, porque el
coste depende del tamaño del banco y de si se calculan alineamientos. La segunda es
que \cod{load\_goa\_annotations} tarda más de dos horas y es I/O contra la base:
\textbf{no se paraleliza entre máquinas}, porque una sola máquina posee el estado.

\subsection{Proyección de la rejilla cruzada}

La rejilla cruzada de sustrato por política de banco son $22 \times 16 = 352$ arms:
las 25 configuraciones menos las tres capas de \cod{ankh-base} que no son evaluables.
Con la comparación en cascada y no todos contra todos:

\begin{center}
\begin{tabularx}{\linewidth}{@{}X r r@{}}
\toprule
partida & unidades & coste\\
\midrule
codificación pendiente                   & 11 configuraciones & 120 h GPU\\
predicción                               & 352 arms          & 401 h CPU\\
evaluación, 7 cortes por 2 variantes     & 4{,}928 corridas    & 238 h CPU\\
comparación en cascada                   & unas 4{,}300        & 52 h CPU\\
\midrule
\textbf{total}                           &                   & \textbf{811 h}\\
\midrule
la misma rejilla con todos contra todos  & 61{,}776 comparaciones & +618 h, y no válida\\
\bottomrule
\end{tabularx}
\end{center}

Traducido a plazos, suponiendo cuatro trabajos concurrentes por máquina:

\begin{center}
\begin{tabularx}{\linewidth}{@{}X r r@{}}
\toprule
configuración & concurrencia & plazo\\
\midrule
un trabajador en serie            &  1 & 34 días\\
una máquina, cuatro trabajadores  &  4 & 8.5 días\\
cinco máquinas                    & 20 & \textbf{1.7 días}\\
\bottomrule
\end{tabularx}
\end{center}

\subsection{Lo que cinco máquinas desbloquean, en orden de valor}

\begin{enumerate}[leftmargin=1.5em,itemsep=0.4em]
  \item \textbf{Igualar los bancos de donantes.} Es el bloqueo que hoy impide
        cualquier veredicto de sustrato. Once configuraciones tienen 60{,}797
        embeddings frente a 528{,}294, de modo que hay que calcular unos 467{,}500
        embeddings por configuración, y son \emph{once} configuraciones. Al ritmo
        medido, entre 41{,}350 y 43{,}272 embeddings por hora, son unas 11 h de GPU
        por configuración y unas \textbf{120 h en total}, o unas \textbf{24 h sobre
        cinco GPU}. \textbf{Sigue siendo el mejor retorno del informe: desbloquea un
        eje entero por un día de cómputo sobre cinco GPU.}
  \item \textbf{La rejilla cruzada.} 1.7 días frente a 8.5. Es lo que permite
        observar interacciones entre nodos.
  \item \textbf{El encaminamiento con ajuste cruzado.} Multiplica el coste de
        puntuación por el número de pliegues. Con cinco máquinas es rutina; con una
        es el paso que se queda sin hacer, y sin él la canalización no es publicable.
  \item \textbf{La escalera de recencia del banco.} Trece releases de GOA desde la
        160 hasta la 220, de las que hay dos. \textbf{Esta no se beneficia de las
        cinco máquinas}: son unos 20.3~GB de datos y once cargas de más de dos horas
        contra la única base que posee el estado, necesariamente en serie. Son unas
        24 h de reloj y hay que planificarlas como tales.
\end{enumerate}

\begin{aviso}{Un límite que el cómputo no levanta}
Una máquina posee todo el estado: base de datos, cola y almacén de objetos. Las
otras calculan y publican de vuelta. Todo lo que sea ingesta, sellado o migración
es serie por construcción y no mejora con más nodos. Las cinco máquinas aceleran
codificación, predicción, evaluación y comparación, que son el 98\,\% del coste, y
no aceleran el 2\,\% restante. También hay una condición: los dos árboles de código
tienen que estar en la misma revisión, porque código desparejado etiqueta mal los
resultados y reporta éxito.
\end{aviso}

\section{Procedimiento}
\label{sec:procedimiento}

Cada etapa tiene una puerta que puede fallar, y en cada caso se nombra la entrada
sobre la que falla. Una puerta que no puede fallar no es una puerta.

\subsection{Cambios de instrumentación, en orden}

Los once cambios siguen todos la plantilla de tres partes de la
sección~\ref{sec:proposito}: una fila, un guardia que se niega, y dos pruebas.

\begin{center}
\small
\begin{xltabular}{\linewidth}{@{}c L{5.4cm} L{2.2cm} X@{}}
\toprule
\# & cambio & coste & se niega a\\
\midrule
\endfirsthead
\toprule
\# & cambio & coste & se niega a\\
\midrule
\endhead
1  & filtrar por estado la lectura de umbrales & 1 línea & que una declaración retirada gobierne algo\\
2  & publicar el umbral y el resultado de la comparación junto al nodo & pequeño & nada. Hace visible una negativa que hoy se descarta\\
3  & el marco se \emph{designa}, no es la fila más reciente & pequeño & un marco elegido por orden de inserción\\
4  & promover el veredicto emparejado a una tabla tipada & migración y backfill de 9{,}138 lecturas & que borrar un trabajo borre la única evidencia\\
5  & transportar la sexta palabra del vocabulario & pequeño & que un nulo con poder se lea como una pregunta no hecha\\
6  & la declaración como espina de atribución, con clave foránea desde el resultado & migración y 315 filas & un número publicado sin declaración detrás\\
7  & materializar \cod{th\_step} y el estadístico, y añadirlos al sello & migración y re-sellado & comparar dos rejillas distintas, o dos estadísticos distintos\\
8  & retirar o restringir la etiqueta de ventana del sello & re-sellado & que una etiqueta del arnés parta un marco real\\
9  & negarse a comparar arms con bancos de donantes desiguales & un campo y un guardia & el confundido de la sección~\ref{sec:parametros}\\
10 & hacer que los bloqueos estructurales no se lean como evidencia & pequeño & reportar un veredicto que ningún experimento puede mover\\
11 & ligar la ablación de rasgos al esquema versionado & preparado & una ablación que retire familias en lugar de columnas\\
\bottomrule
\caption{Los once cambios de instrumentación. Todos se ejecutan en la máquina que
posee el estado.}
\end{xltabular}
\end{center}

\subsection{Etapas y puertas}

\begin{description}[leftmargin=0pt,itemsep=0.5em,font=\normalfont\bfseries]
  \item[S0. Designar el marco.] Productos: la ventana de selección designada, la
        variante designada y el punto de validación interna. Puerta \textbf{G0}: la
        ventana designada resuelve a una release menor o igual que 227, y la
        comprobación de comparabilidad deja de negarse. \emph{Falla sobre} una
        designación por encima de 227, una ventana cuyo annotation set no existe, o
        una designación que nombre dos evaluation sets a la vez.

  \item[S1. Sacar la evidencia del registro de eventos.] Puerta \textbf{G1}: en una
        transacción que se revierte, intentar borrar un trabajo de comparación debe
        \emph{levantar} por la restricción de la tabla de veredictos, y el recuento
        restaurado debe reproducir las cifras medidas. \textbf{Nada más corre antes
        de G1}, porque toda etapa posterior añade trabajos cuyo borrado se llevaría
        la evidencia.

  \item[S2. Re-sellar bajo los campos corregidos.] Puerta \textbf{G2}: ningún sello
        contiene dos rejillas de umbral ni dos estadísticos. \emph{Falla sobre el
        registro actual}, donde los cuatro sellos contienen las dos rejillas.

  \item[S3. El sustrato, sobre bancos iguales.] Puerta \textbf{G3}: la comparación
        se niega si los bancos difieren, y la negativa es visible en la respuesta.
        \emph{Falla sobre} el censo actual de embeddings.

  \item[S4. El corte de evaluación.] Puerta \textbf{G4}: retirar los resultados con
        corte cambia el número de niveles de algún nodo. \emph{Falla hoy}, porque
        ningún nodo lee el campo del corte.

  \item[S5. Canalización con ajuste cruzado.] Puerta \textbf{G5}: el veredicto
        publicado es fuera de pliegue, y la comparación se niega cuando la partición
        de ajuste y la de evaluación se cortan.

  \item[S6. La medida de validación, una sola vez.] Precondición: todas las puertas
        anteriores superadas y el campeón congelado. Puerta \textbf{G6}: antes de
        calcular nada, el guardia de fuga debe \emph{demostrar una negativa} sobre
        ese marco y después permitir el paso al campeón.
\end{description}

\begin{aviso}{Por qué G6 es la puerta que sostiene todo el protocolo}
Sobre la ventana adoptada, $220 \rightarrow 227$, el guardia de fuga \emph{no puede
levantar}: las tres codificaciones ajustadas que existen declaran corte en la
release 220, y la condición de negativa exige que el corte caiga estrictamente
dentro de la ventana. El guardia está vivo, se invoca, y de forma demostrable no dispara
para ninguna fila de esta base sobre esta ventana. Su única negativa demostrada está
en una prueba sobre otra ventana. La mitad \emph{actúa} de ese guardia queda sin
evidenciar hasta S6, y ahí es donde tiene que demostrarse.
\end{aviso}

\subsection{Reparto entre máquinas}

Una máquina posee el estado: base de datos, cola de trabajos, almacén de objetos, la
interfaz y todo artefacto que la campaña produce. Las demás son nodos de cómputo sin
estado: se suman a la cola, calculan y publican de vuelta. Un nodo de cómputo puede
reiniciarse sin aviso y eso es inocuo por diseño, porque el estado y la cola están en
la otra máquina y una corrida continúa a menor rendimiento hasta que el nodo vuelve.

\begin{center}
\begin{tabularx}{\linewidth}{@{}X l@{}}
\toprule
trabajo & máquina\\
\midrule
codificación (GPU) & nodo de cómputo\\
predicción y recuperación KNN (CPU) & nodo de cómputo\\
evaluación y estratificación & cualquiera\\
comparación emparejada & la que posee el estado\\
ingesta de GOA y ontologías, migraciones, sellado & la que posee el estado, en serie\\
\bottomrule
\end{tabularx}
\end{center}

\section{Límites: lo que la campaña no va a poder afirmar}
\label{sec:limites}

Esta sección existe para que ninguna de estas cosas se descubra al final. Cada una
es una afirmación que el diseño \emph{no} soporta, con la razón.

\subsection{Sobre comparaciones externas}

\textbf{Ninguna comparación con las cifras publicadas previamente por el proyecto.}
Son cantidades distintas en tres sentidos a la vez: en espacio de métrica, en
ventana temporal y en arnés de ejecución. Varias de ellas son promedios sobre los
nueve paneles o sobre categorías combinadas, que esta campaña no calcula por
construcción. Ponerlas en la misma tabla sería un error de categoría, no una
comparación desfavorable.

\textbf{Ninguna afirmación frente a métodos basados en alineamiento.} No existe hoy
un comparador de esa clase dentro de la plataforma: el único código que evalúa una
herramienta externa no está registrado como operación y su salida no puede llegar a
la tabla de resultados. La superficie de afirmación viva es \emph{interna y del lado
de la recuperación}: puede ordenar configuraciones entre sí y todavía no puede
situar ninguna frente a un método externo. Eso es exactamente lo que el despliegue
en LAFA resuelve, y es la razón por la que el despliegue es parte del diseño y no un
apéndice.

\subsection{Sobre los nodos bloqueados}

Seis de los diez nodos están bloqueados, cada uno por un cero o por un uno con
nombre, y los ceros no son uno por nodo: no hay
anotaciones de dominios, no hay mapeo de dominios a términos, no hay configuraciones
de puntuación registradas, no hay modelos de reordenación, no hay conjuntos de
datos exportados y no hay coocurrencia de términos. Que estén bloqueados no es un
fallo de la campaña: es el estado del registro, y avanzar uno por uno es el plan.

Lo que sí hay que corregir es que cuatro de esos seis bloqueos son literales en el
código y no lecturas del registro, de modo que imprimirían la misma palabra con un
millón de filas presentes. Un bloqueo estructural no debe poder leerse como
evidencia. Es el cambio 10 de la sección~\ref{sec:procedimiento}.

\subsection{Sobre reproducibilidad}

Lo medido es que la \emph{puntuación} es determinista: 117 de 117 métricas se
reproducen exactamente. Eso es más débil que decir que una corrida reproduce
modelos idénticos, porque hoy no hay modelos. La afirmación correcta es sobre el
scoring, y no debe ampliarse.

\subsection{Sobre lo que no se ha establecido}

\begin{itemize}[leftmargin=1.4em,itemsep=0.3em]
  \item \textbf{Las dos variantes de la ventana difieren} en un factor de 2.45 en
        uno de los paneles y son idénticas en otro. El mecanismo \emph{no está
        establecido}. Cualquier veredicto que las mezcle es nulo, y cualquiera que
        se publique sobre una tiene que nombrar cuál.
  \item \textbf{La ventana retira 426{,}385 anotaciones sobre 47{,}455 proteínas},
        el doble de proteínas de las que ganan anotación, y nadie ha desglosado esa
        cifra por causa: término obsoleto, cambio de relación, o retracción real.
        Hasta que se desglose, la confianza en la ventana está acotada por un número
        sin examinar mayor que las ganancias que la campaña mide.
  \item \textbf{Tres configuraciones no son evaluables}: tres capas intermedias de
        uno de los modelos producen puntuación constante. No deben entrar en ninguna
        rejilla.
  \item \textbf{La fracción de columnas vacías del vector de rasgos no está
        establecida con precisión.} Dos muestreos sucesivos discrepan. Importa para
        la ablación: retirar una columna que es nula en la mayoría de las filas
        devuelve \emph{no contribuye} por una razón que no tiene nada que ver con la
        señal.
\end{itemize}

\subsection{Y una propiedad del diseño que conviene no perder}

El holdout se protege por ausencia y no por disciplina: las releases que no deben
mirarse no están cargadas. Cualquier decisión futura que cargue una de ellas por
comodidad convierte una garantía estructural en una promesa, y las promesas de esta
clase ya se han incumplido en este proyecto sin que nadie lo notase. La fricción es
la característica, no el defecto.


\end{document}
